\begin{abstract}
High-entropy alloys (HEAs) combine high strength, structural stability, and---in carefully chosen compositions---passive-film corrosion resistance that can exceed conventional marine alloys in chloride media. Additive manufacturing (AM) is a natural processing partner for these alloys: rapid solidification suppresses detrimental elemental segregation, refines the microstructure, and can improve pitting resistance relative to cast material, while powder- and wire-based routes enable complex wetted components, large structures, and in-service repair that are uneconomical by casting or forging. This review synthesizes the AM processing of HEAs (LPBF, SEBM, DED/LMD, WAAM, and binder jetting), the as-built and heat-treated microstructures that result, and their electrochemical and corrosion performance in saline and seawater environments, before mapping these capabilities onto concrete marine applications: seawater piping, pumps and valves, heat exchangers, propulsion components, and cladding-based repair of shafts and propellers. We show that the main scientific gaps are environment-specific: long-term natural seawater exposure, biofilm-coupled corrosion, corrosion fatigue and stress corrosion cracking data, and galvanic compatibility of HEA deposits on marine steels are all poorly quantified. The main engineering gaps are feedstock cost, defect control, and the absence of classification-society qualification pathways for HEA compositions. We close with a staged roadmap linking alloy design, process control, and standards-based qualification for marine deployment.
\end{abstract}

\begin{IEEEkeywords}
High-entropy alloys, additive manufacturing, laser powder bed fusion, wire arc additive manufacturing, directed energy deposition, seawater corrosion, pitting, passivation, marine engineering, qualification.
\end{IEEEkeywords}
